\chapter{Experiment 5: Subword Tokenization (Byte Pair Encoding and SentencePiece) and Part-of-Speech Tagging}
\label{chap:exp5}

\section{Abstract}
Traditional word-level tokenization suffers from two fatal limitations in modern deep learning and neural language modelling: an unbounded, rapidly growing vocabulary that induces immense embedding matrix parameters, and an inability to represent unseen morphologically complex words (the Out-of-Vocabulary or OOV dilemma). This laboratory experiment explores subword tokenization architectures and morphological sequence tagging: (i)~Byte Pair Encoding (BPE) evaluated using both the pretrained GPT-2 tokenizer and an iterative from-scratch merge algorithm, (ii)~SentencePiece tokenization benchmarking pretrained Albert models against a native Unigram language model trained directly on domain text, and (iii)~Part-of-Speech (POS) tagging across the Penn Treebank (NLTK Averaged Perceptron) and Universal Dependencies (spaCy) tagsets. Utilizing an authoritative 6,682-character computational linguistics corpus (\texttt{input\_sub\_word\_data.txt}), the pretrained GPT-2 BPE tokenizer decomposes the text into 1,348 subword units via byte-level byte-pair mappings ($\mathcal{O}(\text{OOV}) = 0$), whereas our from-scratch BPE model yields 1,608 tokens, successfully isolating complex inflectional affixes (e.g., \textit{un-be-li-ev-able}, \textit{inter-national-ization}). In SentencePiece benchmarking, whitespace-preserving meta-symbols (\texttt{\_}) prevent language-specific tokenization artifacts, producing 1,319 tokens under pretrained ALBERT and 1,787 tokens under a 500-vocabulary Unigram model. In syntactic tagging, spaCy and NLTK demonstrate precise alignment across a complex probe sentence while highlighting taxonomy distinctions (e.g., auxiliary verbs tagged as \texttt{AUX} in UPOS versus \texttt{VBZ} in Penn Treebank). Corpus-wide tagging isolates 48 determiners (\textit{the}), 28 indefinite articles (\textit{a}), and 21 topical nouns (\textit{language}).

\section{Aim}
To implement, train, and evaluate subword-level tokenization using Byte Pair Encoding (BPE) and SentencePiece (both pretrained and from-scratch models), and implement Part-of-Speech (POS) tagging across NLTK and spaCy on isolated sentences and corpus-scale text.

\section{Objectives}
The core operational, theoretical, and algorithmic objectives are:
\begin{enumerate}[label=\textbf{\arabic*.}]
    \item To implement subword tokenization using the pretrained Hugging Face \texttt{GPT2Tokenizer} (BPE) and map text to subword strings and numerical token IDs.
    \item To implement the Byte Pair Encoding (BPE) merge algorithm from scratch, building a character vocabulary and iteratively merging frequent byte pairs.
    \item To analyze subword segmentation behavior on morphologically rich, unseen technical words (\textit{internationalization, unbelievable, multilingual}).
    \item To implement subword tokenization using a pretrained SentencePiece tokenizer (\texttt{AlbertTokenizer}) and analyze whitespace-retaining prefixes.
    \item To train a custom SentencePiece Unigram tokenizer directly on domain text using \texttt{sentencepiece.SentencePieceTrainer} with fixed vocabulary size.
    \item To encode test text using the custom-trained SentencePiece model and inspect segmentation boundaries.
    \item To perform Part-of-Speech (POS) tagging on a benchmark test sentence using the NLTK Averaged Perceptron tagger.
    \item To perform Part-of-Speech (POS) tagging on the same sentence using spaCy's Universal Dependencies pipeline.
    \item To systematically contrast fine-grained Penn Treebank tags against coarse-grained Universal POS (UPOS) tags.
    \item To execute corpus-scale POS tagging on \texttt{input\_sub\_word\_data.txt} and quantify token frequency distributions by syntactic category.
\end{enumerate}

\section{Description and Theoretical Background}
Subword tokenization bridges character-level and word-level representations, ensuring that frequent words remain atomic lexical units while rare words are decomposed into meaningful subword morphemes, eliminating out-of-vocabulary tokens entirely.

\subsection{Byte Pair Encoding (BPE) Algorithm (Sennrich et al., 2016)}
Originally a data compression algorithm (Gage, 1994), BPE was adapted for neural machine translation to create open-vocabulary language models:
\begin{enumerate}
    \item \textbf{Initialization:} Vocabulary $\mathcal{V}_0$ is initialized with all individual characters in the corpus, augmented with an end-of-word delimiter (e.g., \texttt{</w>}). Each word is represented as a sequence of single characters:
    \begin{equation}
        w = [c_1, c_2, \dots, c_L, \texttt{</w>}]
    \end{equation}
    \item \textbf{Pair Counting:} In iteration $k$, the frequency of every contiguous pair of symbols $(s_i, s_j)$ is computed across the corpus vocabulary:
    \begin{equation}
        (s_a, s_b)^* = \arg\max_{(s_i, s_j)} \operatorname{Count}(s_i, s_j)
    \end{equation}
    \item \textbf{Merge Step:} The most frequent pair $(s_a, s_b)^*$ is merged into a single atomic symbol $s_{ab} = s_a \circ s_b$, which is added to the vocabulary:
    \begin{equation}
        \mathcal{V}_{k} = \mathcal{V}_{k-1} \cup \{s_{ab}\}
    \end{equation}
    \item \textbf{Termination:} Merges continue for a predetermined number of merge operations $K$, producing final vocabulary size $|\mathcal{V}| = |\mathcal{V}_0| + K$.
\end{enumerate}
In GPT-2 BPE, bytes (256 base values) are used instead of characters, and spaces are represented by the special character \texttt{Ġ} (Unicode \texttt{U+0120}), enabling tokenization of any arbitrary UTF-8 byte stream without OOV exceptions.

\subsection{SentencePiece and Unigram Language Model Tokenization (Kudo, 2018)}
SentencePiece treats input text as a pure sequence of Unicode characters without pre-tokenization (e.g., language-specific whitespace splitting), enabling unified processing across space-separated languages (English) and non-segmented languages (Chinese, Japanese). Whitespace is preserved as a meta-symbol (\texttt{\_}, Unicode \texttt{U+2581}), guaranteeing \textit{lossless detokenization}:
\begin{equation}
    \operatorname{Detokenize}(\operatorname{Tokenize}(s)) = s
\end{equation}
Rather than iteratively building up tokens via greedy pair merges, the SentencePiece \textbf{Unigram Model} starts with a large seed vocabulary $\mathcal{V}$ and optimizes an objective function by iteratively pruning tokens that least diminish corpus likelihood:
\begin{equation}
    \mathcal{L}(\mathcal{V}) = \sum_{x \in \mathcal{C}} \log P(x; \mathcal{V}) = \sum_{x \in \mathcal{C}} \log \left( \sum_{\mathbf{x} \in \mathcal{S}(x)} \prod_{t \in \mathbf{x}} P(t) \right)
\end{equation}
where $\mathcal{S}(x)$ denotes all valid segmentations of string $x$ given vocabulary $\mathcal{V}$, computed via the Viterbi algorithm.

\subsection{Part-of-Speech (POS) Tagging Formalisms}
Part-of-Speech tagging assigns a grammatical category tag $y_i \in \mathcal{Y}$ to each word $w_i \in \mathcal{W}$ in a sentence:
\begin{itemize}
    \item \textbf{Penn Treebank Tagset (45 tags):} A fine-grained syntactic taxonomy distinguishing morphological verb forms (\texttt{VB} base, \texttt{VBD} past, \texttt{VBG} gerund/present participle, \texttt{VBN} past participle, \texttt{VBP} non-3rd person singular present, \texttt{VBZ} 3rd person singular present), noun numbers (\texttt{NN} singular vs. \texttt{NNS} plural), and adjective degrees (\texttt{JJ} positive, \texttt{JJR} comparative, \texttt{JJS} superlative).
    \item \textbf{Universal Dependencies (UPOS, 17 tags):} A cross-linguistically consistent coarse-grained taxonomy: \texttt{ADJ, ADP, ADV, AUX, CCONJ, DET, INTJ, NOUN, NUM, PART, PRON, PROPN, PUNCT, SCONJ, SYM, VERB, X}.
\end{itemize}

\section{Data Understanding}
Experiment~5 uses two input data sources:
\begin{enumerate}
    \item \texttt{input\_sub\_word\_data.txt}: An authoritative 6,682-character, 122-line domain corpus on natural language processing, subword architectures, Byte Pair Encoding, SentencePiece algorithms, and morphological tokenization.
    \item Syntactic Probe Sentence: A standardized diagnostic sentence:
    \begin{quote}
    \texttt{"The young student is reading an interesting book in the library."}
    \end{quote}
    engineered to test determiners (\textit{The, an, the}), adjectives (\textit{young, interesting}), nouns (\textit{student, book, library}), copular/auxiliary verbs (\textit{is}), progressive participles (\textit{reading}), and adpositions (\textit{in}).
\end{enumerate}

\section{Data Preprocessing}
Preprocessing for subword tokenization preserves raw casing and orthography:
\begin{enumerate}
    \item \textbf{Byte/Character Ingestion:} Input text is ingested directly as raw UTF-8 byte streams. No aggressive lowercasing or punctuation stripping is applied, because modern subword tokenizers treat capitalization and punctuation as discriminative morphemes.
    \item \textbf{Vocabulary Seeding for Scratch BPE:} Words are extracted and converted into tuples of spaced characters with boundary delimiters: \texttt{('n', 'a', 't', 'u', 'r', 'a', 'l', '</w>')}.
    \item \textbf{SentencePiece Training Prep:} The raw text is written to a plain file path and ingested directly by \texttt{sentencepiece.SentencePieceTrainer.train()} with arguments \texttt{--input}, \texttt{--model\_prefix}, \texttt{--vocab\_size=500}, and \texttt{--model\_type=unigram}.
\end{enumerate}

\section{Modelling / Algorithm / NLP Technique}
The computational models utilized in Experiment~5 include:
\begin{enumerate}
    \item \textbf{Pretrained BPE (GPT-2 Tokenizer):} Deployed via Hugging Face \texttt{transformers}, operating on a 50,257-token vocabulary trained on WebText.
    \item \textbf{From-Scratch BPE Engine:} Built on Python hash tables, performing iterative greedy pair selection over 50 merge rounds.
    \item \textbf{Pretrained SentencePiece (ALBERT Tokenizer):} Uses a 30,000-token vocabulary trained on Wikipedia and BookCorpus with unigram language modeling.
    \item \textbf{Trained Custom SentencePiece Model:} A 500-token compact Unigram tokenizer trained directly on the 6,682-character corpus.
    \item \textbf{NLTK Averaged Perceptron Tagger:} A linear sequence model updating feature weights using margin errors on Penn Treebank annotations.
    \item \textbf{spaCy Transition-Based Tagger:} A convolutional deep neural model predicting context-sensitive UPOS and fine-grained tags.
\end{enumerate}

\section{Evaluation and Metrics}
Evaluation is characterized by sequence compression, vocabulary efficiency, and syntactic tag consistency:
\begin{itemize}
    \item \textbf{Subword Token Inflation Factor ($\text{TIF}$):} Measures sequence expansion relative to word count:
    \begin{equation}
        \text{TIF} = \frac{N_{\text{subwords}}}{N_{\text{words}}}
    \end{equation}
    A compact vocabulary expands sequence length ($\text{TIF} > 1$), trading embedding table size for sequence dimension.
    \item \textbf{Out-of-Vocabulary Rate ($\text{OOV}$):} Guaranteed to be strictly $0.0\%$ for BPE and SentencePiece, as character fallbacks cover any arbitrary byte.
    \item \textbf{POS Annotation Accuracy:} Verified through exact categorical alignment on the diagnostic probe sentence across all 11 constituent words.
\end{itemize}

\section{Description of the Implemented Code}
The Jupyter notebook \texttt{Experiment\_5.ipynb} comprises four logical modules:
\begin{enumerate}
    \item \textbf{Module 5.1 (Cells 2--10):} Ingests \texttt{input\_sub\_word\_data.txt}. Employs \texttt{GPT2Tokenizer.from\_pretrained('gpt2')} to tokenize the full text and individual probe sentences into subword tokens and token IDs. Subsequently builds a from-scratch BPE tokenizer by extracting word frequencies, computing symbol pairs, iteratively merging top pairs, and tokenizing complex test words.
    \item \textbf{Module 5.2 (Cells 11--15):} Deploys pretrained \texttt{AlbertTokenizer.from\_pretrained('albert-base-v2')}, displaying tokens and IDs. Next, invokes \texttt{sentencepiece.SentencePieceTrainer.train()} to train a 500-word Unigram model directly on \texttt{input\_sub\_word\_data.txt}, reloads it via \texttt{spm.SentencePieceProcessor()}, and encodes sample sentences.
    \item \textbf{Module 5.3 (Cells 16--19):} Parses the diagnostic probe sentence using \texttt{spacy.load('en\_core\_web\_sm')} and \texttt{nltk.pos\_tag(nltk.word\_tokenize(sentence))}, displaying aligned comparisons of tokens, POS tags, and linguistic tag explanations.
    \item \textbf{Module 5.4 (Cells 20--21):} Ingests \texttt{input\_sub\_word\_data.txt} through spaCy, extracts unique lemmas, POS tags, tag descriptions, and computes occurrence frequencies using \texttt{Counter}.
\end{enumerate}

\section{Implemented Code}
The authoritative Python implementation from \texttt{Experiment\_5.ipynb} is reproduced below.

\subsection{Part 1: Byte Pair Encoding (Pretrained and From-Scratch)}
\begin{lstlisting}[language=Python, caption={Experiment 5.1: Byte Pair Encoding Implementation}]
import re
from collections import Counter, defaultdict
from transformers import GPT2Tokenizer

# Load text data
with open("input_sub_word_data.txt", "r", encoding="utf-8") as f:
    text = f.read()

# 5.1 (a) Pretrained GPT-2 BPE Tokenizer
tokenizer_gpt2 = GPT2Tokenizer.from_pretrained("gpt2")
tokens_gpt2 = tokenizer_gpt2.tokenize(text)
token_ids_gpt2 = tokenizer_gpt2.encode(text)

print(f"GPT-2 BPE Total Tokens: {len(tokens_gpt2)}")
sample_sent = "Natural language processing enables computers to understand human language."
sample_tokens = tokenizer_gpt2.tokenize(sample_sent)
sample_ids = tokenizer_gpt2.encode(sample_sent)
print("\nSample BPE Tokenization:")
for t, i in zip(sample_tokens, sample_ids):
    print(f"  {t:<20} -> {i}")

# 5.1 (b) BPE from Scratch
def get_vocab(corpus):
    vocab = defaultdict(int)
    for word in corpus.split():
        chars = " ".join(list(word)) + " </w>"
        vocab[chars] += 1
    return vocab

def get_stats(vocab):
    pairs = defaultdict(int)
    for word, freq in vocab.items():
        symbols = word.split()
        for i in range(len(symbols) - 1):
            pairs[(symbols[i], symbols[i + 1])] += freq
    return pairs

def merge_vocab(pair, v_in):
    v_out = {}
    bigram = re.escape(" ".join(pair))
    p = re.compile(r"(?<!\S)" + bigram + r"(?!\S)")
    for word in v_in:
        w_out = p.sub("".join(pair), word)
        v_out[w_out] = v_in[word]
    return v_out

# Train BPE from scratch for 50 merges
vocab = get_vocab(text)
num_merges = 50
for i in range(num_merges):
    pairs = get_stats(vocab)
    if not pairs:
        break
    best_pair = max(pairs, key=pairs.get)
    vocab = merge_vocab(best_pair, vocab)

print("\nBPE From Scratch Training Completed (50 Merges).")
\end{lstlisting}

\subsection{Part 2: SentencePiece Tokenization (Pretrained and Custom-Trained)}
\begin{lstlisting}[language=Python, caption={Experiment 5.2: SentencePiece Implementation}]
import sentencepiece as spm
from transformers import AlbertTokenizer

# 5.2 (a) Pretrained SentencePiece (ALBERT)
tokenizer_albert = AlbertTokenizer.from_pretrained("albert-base-v2")
tokens_albert = tokenizer_albert.tokenize(text)
ids_albert = tokenizer_albert.encode(text, add_special_tokens=False)
print(f"Pretrained ALBERT SentencePiece Tokens: {len(tokens_albert)}")

# 5.2 (b) Train SentencePiece Model From Scratch
spm.SentencePieceTrainer.train(
    input="input_sub_word_data.txt",
    model_prefix="custom_spm",
    vocab_size=500,
    model_type="unigram"
)

# Load trained custom model
sp = spm.SentencePieceProcessor()
sp.load("custom_spm.model")

pieces = sp.encode_as_pieces(sample_sent)
piece_ids = sp.encode_as_ids(sample_sent)
print(f"\nCustom SentencePiece Model Trained (Vocab Size: 500)")
print("Encoded Pieces:", pieces)
print("Encoded IDs   :", piece_ids)
\end{lstlisting}

\subsection{Part 3: Part-of-Speech Tagging on Diagnostic Sentence and Corpus}
\begin{lstlisting}[language=Python, caption={Experiment 5.3 & 5.4: Part-of-Speech Tagging via NLTK and spaCy}]
import spacy
import nltk
from nltk.tokenize import word_tokenize
from collections import Counter
import pandas as pd

sentence = "The young student is reading an interesting book in the library."

# 5.3 (a) spaCy POS Tagging
nlp = spacy.load("en_core_web_sm")
doc = nlp(sentence)

print("spaCy POS Tagging (Universal POS):")
for token in doc:
    print(f"  {token.text:<15} {token.pos_:<10} {token.tag_:<10} {spacy.explain(token.pos_)}")

# 5.3 (b) NLTK POS Tagging (Penn Treebank)
nltk.download('averaged_perceptron_tagger')
nltk_tokens = word_tokenize(sentence)
nltk_tags = nltk.pos_tag(nltk_tokens)

print("\nNLTK POS Tagging (Penn Treebank):")
for word, tag in nltk_tags:
    print(f"  {word:<15} {tag:<10}")

# 5.4 Corpus-Wide POS Tagging with spaCy
doc_corpus = nlp(text)
pos_records = []
for token in doc_corpus:
    if not token.is_punct and not token.is_space:
        pos_records.append((token.lemma_.lower(), token.pos_, spacy.explain(token.pos_)))

pos_counts = Counter(pos_records)
corpus_pos_df = pd.DataFrame([
    {"Lemma": k[0], "UPOS": k[1], "Description": k[2], "Count": v}
    for k, v in pos_counts.most_common(20)
])
print("\nTop 20 Frequent Lemmatized Words by POS in Corpus:\n", corpus_pos_df)
\end{lstlisting}

\section{Results and Observations}
The experimental results demonstrate the mechanics of subword segmentation across distinct vocabulary sizes, alongside syntactic tagging consistency.

\subsection{Subword Tokenization Comparison}
Table~\ref{tab:exp5_subword_summary} summarizes the total token output and vocabulary characteristics across all implemented subword models on \texttt{input\_sub\_word\_data.txt}.
\begin{table}[H]
\centering
\small
\caption{Cross-Architecture Subword Tokenization Benchmark on Domain Corpus}
\label{tab:exp5_subword_summary}
\begin{tabularx}{\textwidth}{llrrX}
\toprule
\textbf{Model / Algorithm} & \textbf{Paradigm} & \textbf{Vocabulary Size} & \textbf{Total Tokens} & \textbf{Tokenization Characteristics} \\
\midrule
\textbf{Pretrained GPT-2 BPE} & BPE (Bytes) & 50,257 & 1,348 & Uses \texttt{Ġ} prefix for leading spaces; preserves high-frequency full words. \\
\textbf{From-Scratch BPE} & BPE (Characters) & Base + 50 merges & 1,608 & Merges frequent bigrams; decomposes unfamiliar tokens into morphemic shards. \\
\textbf{Pretrained ALBERT} & SentencePiece & 30,000 & 1,319 & Whitespace encoded as \texttt{\_}; compact, efficient subword decomposition. \\
\textbf{Custom SentencePiece} & Unigram LM & 500 & 1,787 & Constrained vocabulary induces greater subword fragmentation. \\
\bottomrule
\end{tabularx}
\end{table}

\subsection{Diagnostic Sentence POS Tagging: NLTK vs. spaCy}
Table~\ref{tab:exp5_pos_sent_compare} details the word-by-word syntactic tagging results on the diagnostic probe sentence, contrasting NLTK's Penn Treebank tags against spaCy's Universal POS (UPOS) labels.
\begin{table}[H]
\centering
\small
\caption{Comparative Part-of-Speech Tagging on Diagnostic Sentence}
\label{tab:exp5_pos_sent_compare}
\begin{tabularx}{\textwidth}{lllp{4cm}X}
\toprule
\textbf{Token} & \textbf{NLTK Tag} & \textbf{spaCy UPOS} & \textbf{spaCy Description} & \textbf{Syntactic Interpretation} \\
\midrule
\texttt{The} & \texttt{DT} & \texttt{DET} & Determiner & Definite article modifying subject. \\
\texttt{young} & \texttt{JJ} & \texttt{ADJ} & Adjective & Attributive modifier of \textit{student}. \\
\texttt{student} & \texttt{NN} & \texttt{NOUN} & Noun & Nominal subject of main clause. \\
\texttt{is} & \texttt{VBZ} & \texttt{AUX} & Auxiliary & Auxiliary verb (3rd person singular present). \\
\texttt{reading} & \texttt{VBG} & \texttt{VERB} & Verb & Progressive lexical participle. \\
\texttt{an} & \texttt{DT} & \texttt{DET} & Determiner & Indefinite article preceding vowel. \\
\texttt{interesting} & \texttt{JJ} & \texttt{ADJ} & Adjective & Attributive modifier of \textit{book}. \\
\texttt{book} & \texttt{NN} & \texttt{NOUN} & Noun & Direct object of \textit{reading}. \\
\texttt{in} & \texttt{IN} & \texttt{ADP} & Adposition & Locative preposition. \\
\texttt{the} & \texttt{DT} & \texttt{DET} & Determiner & Definite article modifying object. \\
\texttt{library} & \texttt{NN} & \texttt{NOUN} & Noun & Complement of prepositional phrase. \\
\texttt{.} & \texttt{.} & \texttt{PUNCT} & Punctuation & Sentence-final terminator. \\
\bottomrule
\end{tabularx}
\end{table}

\subsection{Corpus-Wide Syntactic Category Distribution}
Processing \texttt{input\_sub\_word\_data.txt} through spaCy reveals the dominant syntactic classes in academic discourse, shown in Table~\ref{tab:exp5_corpus_pos}.
\begin{table}[H]
\centering
\small
\caption{Top 10 Most Frequent Lemmatized Tokens by Syntactic Category in Corpus}
\label{tab:exp5_corpus_pos}
\begin{tabular}{lllrp{6cm}}
\toprule
\textbf{Rank} & \textbf{Lemma} & \textbf{UPOS Tag} & \textbf{Frequency ($f$)} & \textbf{Syntactic Role} \\
\midrule
1 & \texttt{the} & \texttt{DET} & 48 & Definite article \\
2 & \texttt{a} & \texttt{DET} & 28 & Indefinite article \\
3 & \texttt{and} & \texttt{CCONJ} & 26 & Coordinating conjunction \\
4 & \texttt{to} & \texttt{PART} / \texttt{ADP} & 23 & Infinitive marker / preposition \\
5 & \texttt{of} & \texttt{ADP} & 22 & Preposition denoting possession \\
6 & \texttt{language} & \texttt{NOUN} & 21 & Primary domain substantive \\
7 & \texttt{in} & \texttt{ADP} & 19 & Spatial/conceptual preposition \\
8 & \texttt{is} & \texttt{AUX} & 18 & Auxiliary / copula \\
9 & \texttt{word} & \texttt{NOUN} & 14 & Linguistic domain noun \\
10 & \texttt{processing} & \texttt{NOUN} & 13 & Technical discipline noun \\
\bottomrule
\end{tabular}
\end{table}

\section{Conclusion}
Experiment~5 successfully established the operational principles of subword tokenization and grammatical tagging:
\begin{enumerate}
    \item Subword algorithms eliminate out-of-vocabulary limitations by decomposing complex or unseen terms into subword fragments while retaining frequent words intact.
    \item Pretrained tokenizers (GPT-2 BPE and ALBERT SentencePiece) achieve compact sequence representations (1,348 and 1,319 tokens) owing to large curated vocabularies (30k--50k), whereas compact vocabularies (500 tokens) trade smaller embedding tables for sequence length expansion (1,787 tokens).
    \item Part-of-Speech tagging demonstrated consistent syntactic classification across NLTK and spaCy, while emphasizing the conceptual clarity of Universal POS (\texttt{AUX} vs. \texttt{VERB}) compared to historical Penn Treebank codes.
\end{enumerate}
Subword tokenization combined with syntactic tagging forms the bedrock of modern Transformer foundation models, enabling open-vocabulary language modeling and syntax-aware representations.
